\documentclass[10pt,journal,compsoc]{IEEEtran}

\usepackage[utf8]{inputenc}
\usepackage{amsmath,amssymb,amsfonts}
\usepackage{graphicx}
\usepackage{booktabs}
\usepackage{multirow}
\usepackage{listings}
\usepackage{xcolor}
\usepackage{hyperref}
\usepackage{subcaption}
\usepackage{algorithm}
\usepackage{algpseudocode}
\usepackage{cite}

% ====================================================================
% UNIFIED TYPOGRAPHY SPECIFICATION: EXACTLY ONE SINGLE FONT FAMILY
% We enforce Times Roman (ptm / mathptmx) across the entire paper:
% body text, headings, math formulas, captions, and code listings.
% ====================================================================
\usepackage{mathptmx}
\renewcommand{\familydefault}{\rmdefault}
\renewcommand{\sfdefault}{ptm}
\renewcommand{\ttdefault}{ptm}

% Code listing styling enforcing the exact same Times Roman typeface
\definecolor{codegreen}{rgb}{0,0.5,0}
\definecolor{codegray}{rgb}{0.4,0.4,0.4}
\definecolor{codepurple}{rgb}{0.5,0,0.5}
\definecolor{backcolour}{rgb}{0.97,0.97,0.97}

\lstdefinestyle{pythonstyle}{
    backgroundcolor=\color{backcolour},   
    commentstyle=\fontfamily{ptm}\selectfont\itshape\color{codegreen},
    keywordstyle=\fontfamily{ptm}\selectfont\bfseries\color{black},
    numberstyle=\tiny\fontfamily{ptm}\selectfont\color{codegray},
    stringstyle=\fontfamily{ptm}\selectfont\color{codepurple},
    basicstyle=\fontfamily{ptm}\selectfont\scriptsize,
    breakatwhitespace=false,         
    breaklines=true,                 
    captionpos=b,                    
    keepspaces=true,                 
    numbers=left,                    
    numbersep=5pt,                  
    showspaces=false,                
    showstringspaces=false,
    showtabs=false,                  
    tabsize=2
}
\lstset{style=pythonstyle}

\hypersetup{
    colorlinks=true,
    linkcolor=black,
    filecolor=black,      
    urlcolor=blue,
    citecolor=black
}

\begin{document}

\title{Can You Beat Me at Chess?\\Design, Training, and Squeeze-and-Excitation Residual Modeling of \texttt{nom\_bot\_chess}}

\author{L\^e~Th\`anh~Nam%
\thanks{L\^e Th\`anh Nam is an independent artificial intelligence researcher and lead developer of \texttt{nom\_bot\_chess} (Email: nam29@nom-documents.org). Code repository and open checkpoints are publicly maintained at: \url{https://github.com/lethanhnam/nom_bot_chess}}%
}

\markboth{IEEE Transactions on Games,~Vol.~X, No.~Y,~November~2026}%
{Nam: Can You Beat Me at Chess? Design and Evaluation of \texttt{nom\_bot\_chess}}

\IEEEtitleabstractindextext{%
\begin{abstract}
Chess has long stood as an intellectual benchmark for artificial intelligence. While superhuman engines such as Stockfish and AlphaZero demonstrate absolute dominance via massive compute and tree search, developing an accessible, intermediate-to-expert level agent (Elo 1500--1800+) that exhibits human-like grandmaster positional intuition while executing strictly within accessible cloud resources remains a compelling challenge. In this work, we introduce \texttt{nom\_bot\_chess}, an end-to-end deep learning chess engine built on an 18-plane spatial board representation and an 8-block, 10.3-million-parameter dual-head Deep Residual Neural Network (ResNet). The model was trained via supervised learning on $956,713$ discrete board states extracted from $25,000$ high-tier competitive games ($2400+$ vs $2200+$ rating) from the Lichess Elite Database (August 2020). Using a strict, game-level partition of $90\%$ training ($860,702$ samples), $5\%$ validation ($47,660$ samples), and $5\%$ independent testing ($48,351$ samples) to eliminate sequential data leakage, our model achieves a top-1 move accuracy of $35.97\%$, top-3 accuracy of $61.92\%$, and top-5 accuracy of $73.69\%$ on unseen positions. Integrated with a policy-ordered minimax search with alpha-beta pruning, the bot demonstrates tactical soundness and secured a decisive 39-move checkmate victory in baseline competitive benchmarking. Furthermore, we provide an extensive theoretical and architectural investigation into Squeeze-and-Excitation Residual Networks (SE-ResNet), formulating channel-wise feature recalibration to dynamically prioritize crucial piece relationships across phases of the game.
\end{abstract}

\begin{IEEEkeywords}
Computer Chess, Deep Residual Learning, Squeeze-and-Excitation Networks, Alpha-Beta Pruning, Supervised Learning, Game-Playing AI, Lichess Elite Database.
\end{IEEEkeywords}}

\maketitle
\IEEEdisplaynontitleabstractindextext
\IEEEpeerreviewmaketitle

\section{Introduction}
\IEEEPARstart{S}{ince} Garry Kasparov conceded his historic match to Deep Blue in 1997, chess has served as humanity's premier sandbox for artificial intelligence. For decades, computer chess was synonymous with handcrafted evaluation polynomials paired with heavily optimized brute-force alpha-beta tree searches. When DeepMind introduced AlphaZero~\cite{silver2018alphazero}, it revealed that deep residual convolutional networks, when coupled with Monte Carlo Tree Search (MCTS), could learn profound strategic concepts without human opening books or handcrafted endgame tables.

Yet, for ordinary chess enthusiasts, students, and independent practitioners, a question naturally arises: \textit{Can you beat me at chess?} Beyond the philosophical duel between human intuition and silicon calculation, there lies a practical engineering question: \textit{Can an individual researcher construct a robust, human-competitive chess AI (Elo 1500--1800+) entirely on free cloud resources (such as Google Colab) using high-grade grandmaster data, without requiring server clusters or weeks of compute?}

When an intermediate human player (Elo 1500) sits across a chessboard, their play is characterized by two distinct faculties:
\begin{enumerate}
    \item \textbf{Strategic Intuition:} Recognizing harmonious pawn structures, open files, piece development, and king safety within a fraction of a second.
    \item \textbf{Tactical Concrete Calculation:} Calculating 2 to 3 plies ahead to ensure a piece is not left hanging (\textit{en prise}) and to exploit immediate tactical opportunities (forks, pins, skewers).
\end{enumerate}

To embody these dual faculties, we conceived and implemented \texttt{nom\_bot\_chess}. Unlike pure heuristic engines that feel mechanical, or pure deep policy networks that occasionally blunder their queen due to a rare probability blindspot, \texttt{nom\_bot\_chess} unites a dual-head Deep Residual Network (trained on nearly one million grandmaster positions) with an efficient policy-ordered alpha-beta search.

The remainder of this paper is organized as follows: Section~II details the origin, scale, and partition methodology of the Lichess Elite Dataset. Section~III presents the step-by-step pipeline implementation, including tensor encoding, action mapping, and network design. Section~IV analyzes empirical training dynamics, independent test evaluations, and game termination. Section~V delivers a comprehensive theoretical, mathematical, and algorithmic investigation of Squeeze-and-Excitation Residual Networks (SE-ResNet) for chess feature modeling. Section~VI describes deployment and real-time interaction, followed by concluding remarks in Section~VII.

\section{Dataset Origin, Scale, and Partitioning}
\label{sec:dataset}

\subsection{Data Source: Lichess Elite Database}
A critical requirement for training an intermediate-to-expert engine via supervised learning is training exclusively on clean, high-mastery moves. Training on general online games (where players rated 1000--1400 commit frequent tactical blunders) contaminates the policy head with suboptimal move choices.

We utilized the **Lichess Elite Database** curated by Tony Rottella, specifically the monthly archive for **August 2020** (\texttt{lichess\_elite\_2020-08.zip}), hosted at \url{https://database.nikonoel.fr/}. This database is subject to rigorous inclusion criteria:
\begin{itemize}
    \item \textbf{Rating Threshold:} Only standard rated games where both players meet high rating thresholds (Player A $\ge 2400$, Player B $\ge 2200$) are retained.
    \item \textbf{Time Controls:} Bullet games are strictly excluded to avoid chaotic time-trouble inaccuracies; only Blitz, Rapid, and Classical games are preserved.
    \item \textbf{Volume:} The raw August 2020 archive contains $378,687$ complete games, consuming $103\text{ MB}$ compressed (approximately $550\text{ MB}$ in raw Portable Game Notation [PGN] format).
\end{itemize}

\subsection{Sampling and Feature Extraction}
To enable rapid, memory-bounded extraction within Google Colab's standard 12~GB virtual memory constraint, we extracted mid-game states from $N_{\text{games}} = 25,000$ games. To prevent redundancy and temporal auto-correlation between identical consecutive board states:
\begin{enumerate}
    \item The opening 6 plies ($3$ full moves) are skipped, as they represent ubiquitous book moves (e.g., 1.e4 e5 2.Nf3 Nc6).
    \item Positions are subsampled at a stride of $\Delta k = 2$ plies.
    \item Games ending in an indeterminate result (\texttt{*}) are filtered out.
\end{enumerate}
This pipeline yielded a total corpus of **$N = 956,713$ discrete board states**.

\subsection{Strict Game-Level Partitioning ($90\% / 5\% / 5\%$)}
A frequent methodological flaw in game AI literature is \textit{position-level random shuffling} before splitting into train, validation, and test sets. In chess, positions from plies $k$ and $k+1$ within the same game share over $95\%$ identical piece configurations. Shuffling across positions leaks identical opening and endgame positions between training and testing, artificially inflating test accuracy.

To maintain strict scientific integrity, we implemented **Game-Level Partitioning** using modulo binning on the unique game index $g \in [0, N_{\text{games}}-1]$:
\begin{equation}
\text{Split}(g) = 
\begin{cases} 
\text{Train}, & \text{if } (g \bmod 20) < 18 \quad (90\%) \\
\text{Validation}, & \text{if } (g \bmod 20) = 18 \quad (5\%) \\
\text{Test}, & \text{if } (g \bmod 20) = 19 \quad (5\%)
\end{cases}
\end{equation}

Table~\ref{tab:split_summary} details the exact distribution of samples across the three partitions.

\begin{table}[htbp]
\centering
\caption{Dataset Partition Distribution ($90/5/5$ Game-Level Split)}
\label{tab:split_summary}
\begin{tabular}{lrrr}
\toprule
\textbf{Partition} & \textbf{Games} & \textbf{Sampled States} & \textbf{Proportion} \\
\midrule
Training Set       & 22,500         & 860,702                 & $90.0\%$ \\
Validation Set     & 1,250          & 47,660                  & $5.0\%$ \\
Independent Test   & 1,250          & 48,351                  & $5.1\%$ \\
\midrule
\textbf{Total}     & \textbf{25,000}& \textbf{956,713}        & \textbf{100.0\%} \\
\bottomrule
\end{tabular}
\end{table}

The extracted tensors were compressed into an uncorrupted serialized array archive (\texttt{dataset\_split\_cache\_25000.npz}) measuring approximately $240\text{ MB}$, allowing subsequent execution without re-parsing raw PGN text.

\section{System Architecture and Implementation}
\label{sec:architecture}

The end-to-end development of \texttt{nom\_bot\_chess} is structured as an eleven-stage procedural pipeline executed within Google Colab under hardware acceleration (Nvidia Tesla T4 GPU, 16~GB VRAM).

\subsection{18-Plane Spatial Board Representation}
Standard piece-square vectors do not preserve 2D spatial locality. Convolutional layers require multi-channel grid representations. We formulate the board as a tensor $\mathbf{X} \in \mathbb{R}^{18 \times 8 \times 8}$ oriented from the canonical perspective of the side to move (the active player).

The 18 feature channels are defined as follows:
\begin{itemize}
    \item \textbf{Planes 0--5 (Active Pieces):} Binary masks indicating the active player's Pawns, Knights, Bishops, Rooks, Queens, and King.
    \item \textbf{Planes 6--11 (Opponent Pieces):} Binary masks indicating the opponent's Pawns, Knights, Bishops, Rooks, Queens, and King.
    \item \textbf{Planes 12--15 (Castling Rights):} Uniform scalar planes for Active Kingside, Active Queenside, Opponent Kingside, and Opponent Queenside castling rights.
    \item \textbf{Plane 16 (En Passant):} One-hot spatial indicator representing the en-passant target square, if available.
    \item \textbf{Plane 17 (Turn Color):} Uniform plane set to $1.0$ if White to move, or $0.0$ if Black to move.
\end{itemize}

\textbf{Concrete 8$\times$8 Spatial Channel Illustrations:}
To provide an immediate, concrete understanding of how raw chess positions are mapped onto sparse binary planes, Fig.~\ref{fig:channel_matrices} illustrates three representative $8 \times 8$ channels directly from this pipeline:
\begin{enumerate}
    \item \textbf{Channel 0 (Active Pawns):} Reflects pawn positions. Following an opening move such as $1.\text{e4}$, the pawn on the e-file advances to rank 4, resulting in a single $1.0$ on square \texttt{e4} and seven $1.0$ entries along rank 2.
    \item \textbf{Channel 1 (Active Knights):} Demonstrates localized piece placement. With the queenside knight residing at \texttt{b1} and the kingside knight developed to \texttt{f3}, the matrix exhibits exactly two non-zero activations.
    \item \textbf{Channel 16 (En Passant Square):} Represents a transient tactical target. When Black moves \texttt{e7-e5}, square \texttt{e6} becomes the singular active en-passant coordinate ($1.0$), with all other 63 squares remaining $0.0$.
\end{enumerate}
In contrast to sparse piece masks, global game variables such as castling rights (Planes 12--15) and active turn color (Plane 17) are broadcast as uniform scalar planes filled entirely with $1.0$ or $0.0$, allowing spatial convolutional filters to cross-reference global state anywhere on the board.

\begin{figure}[!htbp]
\centering
\begin{subfigure}[b]{\linewidth}
\begin{lstlisting}[style=pythonstyle, numbers=none]
Rank 8: [0, 0, 0, 0, 0, 0, 0, 0]
Rank 7: [0, 0, 0, 0, 0, 0, 0, 0]
Rank 6: [0, 0, 0, 0, 0, 0, 0, 0]
Rank 5: [0, 0, 0, 0, 0, 0, 0, 0]
Rank 4: [0, 0, 0, 0, 1, 0, 0, 0] (e4)
Rank 3: [0, 0, 0, 0, 0, 0, 0, 0]
Rank 2: [1, 1, 1, 1, 0, 1, 1, 1]
Rank 1: [0, 0, 0, 0, 0, 0, 0, 0]
         a  b  c  d  e  f  g  h
\end{lstlisting}
\caption{Channel 0: Active Pawns (after 1.e4)}
\end{subfigure}

\vspace{2mm}
\begin{subfigure}[b]{\linewidth}
\begin{lstlisting}[style=pythonstyle, numbers=none]
Rank 8: [0, 0, 0, 0, 0, 0, 0, 0]
Rank 7: [0, 0, 0, 0, 0, 0, 0, 0]
Rank 6: [0, 0, 0, 0, 0, 0, 0, 0]
Rank 5: [0, 0, 0, 0, 0, 0, 0, 0]
Rank 4: [0, 0, 0, 0, 0, 0, 0, 0]
Rank 3: [0, 0, 0, 0, 0, 1, 0, 0] (f3)
Rank 2: [0, 0, 0, 0, 0, 0, 0, 0]
Rank 1: [0, 1, 0, 0, 0, 0, 0, 0] (b1)
         a  b  c  d  e  f  g  h
\end{lstlisting}
\caption{Channel 1: Active Knights (Nf3 \& Nb1)}
\end{subfigure}

\vspace{2mm}
\begin{subfigure}[b]{\linewidth}
\begin{lstlisting}[style=pythonstyle, numbers=none]
Rank 8: [0, 0, 0, 0, 0, 0, 0, 0]
Rank 7: [0, 0, 0, 0, 0, 0, 0, 0]
Rank 6: [0, 0, 0, 0, 1, 0, 0, 0] (e6)
Rank 5: [0, 0, 0, 0, 0, 0, 0, 0]
Rank 4: [0, 0, 0, 0, 0, 0, 0, 0]
Rank 3: [0, 0, 0, 0, 0, 0, 0, 0]
Rank 2: [0, 0, 0, 0, 0, 0, 0, 0]
Rank 1: [0, 0, 0, 0, 0, 0, 0, 0]
         a  b  c  d  e  f  g  h
\end{lstlisting}
\caption{Channel 16: En Passant Target (e6)}
\end{subfigure}
\caption{Spatial $8 \times 8$ matrix visual representation of selected feature channels in \texttt{nom\_bot\_chess}. (a) Active player's pawn positions following $1.\text{e4}$; (b) Active player's knights with a developed knight on \texttt{f3} and an undeveloped knight on \texttt{b1}; (c) One-hot en-passant target square on \texttt{e6} triggered by Black's two-square pawn advance.}
\label{fig:channel_matrices}
\end{figure}

\textbf{Canonical Perspective Mirroring:} When it is Black's turn to move, ranks are inverted ($r' = 7 - r$) and piece colors are transposed. Consequently, the network always sees the active player moving upwards towards rank 7. This cuts the effective state-space in half and enforces tactical symmetry across White and Black.

Listing~\ref{lst:encoding} demonstrates the implementation of the board encoder.

\begin{lstlisting}[language=Python, caption={Canonical 18-Plane Board Tensor Encoding}, label={lst:encoding}]
def encode_board(board: chess.Board) -> np.ndarray:
    planes = np.zeros((18, 8, 8), dtype=np.float32)
    turn = board.turn  # True = White, False = Black
    piece_map = board.piece_map()
    for sq, piece in piece_map.items():
        rank = chess.square_rank(sq)
        file = chess.square_file(sq)
        if not turn:
            rank = 7 - rank  # Canonical perspective flip
        idx = piece.piece_type - 1
        if piece.color == turn:
            planes[idx, rank, file] = 1.0
        else:
            planes[6 + idx, rank, file] = 1.0
    # Castling & En Passant planes omitted for brevity
    planes[17, :, :] = 1.0 if turn else 0.0
    return planes
\end{lstlisting}

\subsection{Discrete Action Space Formulation}
The action space for general chess comprises all legal moves from any origin square to destination square. We map each move $m = (s_{\text{from}}, s_{\text{to}})$ to a discrete classification index $y \in [0, 4095]$:
\begin{equation}
y = s_{\text{from}} \times 64 + s_{\text{to}}
\end{equation}
where $s_{\text{from}}, s_{\text{to}} \in [0, 63]$. Under Black's turn, squares are reflected via \texttt{chess.square\_mirror}. Pawn promotions default to Queen promotions, which represent $>99\%$ of tactical promotions in grandmaster play.

\subsection{Dual-Head Deep Residual Network Architecture}
The backbone of \texttt{nom\_bot\_chess} is an 8-block Residual Convolutional Neural Network containing $10,323,809$ trainable parameters. 

Each residual block $\mathcal{B}_l$ computes:
\begin{equation}
\mathbf{x}_{l+1} = \text{ReLU}\left(\mathbf{x}_l + \mathcal{F}(\mathbf{x}_l, \{\mathbf{W}_{l,1}, \mathbf{W}_{l,2}\})\right)
\end{equation}
where $\mathcal{F}$ represents two sequential $3 \times 3$ convolutions with Batch Normalization (BN) and ReLU activation:
\begin{equation}
\mathcal{F}(\mathbf{x}) = \text{BN}_2(\text{Conv}_{3\times 3}(\text{ReLU}(\text{BN}_1(\text{Conv}_{3\times 3}(\mathbf{x})))))
\end{equation}

The network branches into two output heads:
\begin{enumerate}
    \item \textbf{Policy Head ($\mathbf{p}$):} Employs a $1 \times 1$ convolution reducing 128 channels to 32, followed by a fully-connected layer outputting raw logits $\mathbf{z} \in \mathbb{R}^{4096}$. During training, softmax cross-entropy against the grandmaster's executed move is applied.
    \item \textbf{Value Head ($v$):} Employs a $1 \times 1$ convolution reducing 128 channels to 16, followed by two dense layers ($1024 \to 128 \to 1$) and a hyperbolic tangent activation $\tanh$, bounding the positional evaluation $v \in [-1.0, +1.0]$.
\end{enumerate}

\begin{lstlisting}[language=Python, caption={Dual-Head ChessResNet Architecture in PyTorch}, label={lst:resnet}]
class ChessResNet(nn.Module):
    def __init__(self, in_channels=18, num_blocks=6, filters=128):
        super().__init__()
        self.in_conv = nn.Sequential(
            nn.Conv2d(in_channels, filters, 3, padding=1, bias=False),
            nn.BatchNorm2d(filters),
            nn.ReLU(inplace=True)
        )
        self.res_blocks = nn.ModuleList([ResBlock(filters) for _ in range(num_blocks)])
        # Policy Head
        self.policy_conv = nn.Sequential(
            nn.Conv2d(filters, 32, kernel_size=1, bias=False),
            nn.BatchNorm2d(32), nn.ReLU(inplace=True)
        )
        self.policy_fc = nn.Linear(32 * 8 * 8, 4096)
        # Value Head
        self.value_conv = nn.Sequential(
            nn.Conv2d(filters, 16, kernel_size=1, bias=False),
            nn.BatchNorm2d(16), nn.ReLU(inplace=True)
        )
        self.value_fc = nn.Sequential(
            nn.Linear(16 * 8 * 8, 128), nn.ReLU(inplace=True),
            nn.Linear(128, 1), nn.Tanh()
        )
\end{lstlisting}

\subsection{Optimization and Multi-Task Loss Function}
The unified loss function combines policy cross-entropy with value mean-squared error:
\begin{equation}
\mathcal{L}_{\text{total}} = \mathcal{L}_{\text{policy}}(\mathbf{p}, \mathbf{y}) + \lambda_v \mathcal{L}_{\text{value}}(v, z)
\end{equation}
where $\mathbf{y}$ is the one-hot target move, $z \in \{+1.0, 0.0, -1.0\}$ is the eventual game outcome from the active player's viewpoint, and $\lambda_v = 0.5$ balances the gradient magnitudes.

Training was conducted using the AdamW optimizer ($\beta_1 = 0.9, \beta_2 = 0.999$, weight decay $10^{-4}$) with an initial learning rate $\eta_0 = 10^{-3}$, modulated via a Cosine Annealing learning rate schedule over 8 epochs with mini-batch size $B = 256$.

\section{Empirical Evaluation and Results}
\label{sec:results}

\subsection{Training Dynamics and Model Selection}
Table~\ref{tab:training_dynamics} tracks the loss and accuracy metrics across all 8 training epochs on the $860,702$ training positions and $47,660$ validation positions.

\begin{table}[htbp]
\centering
\caption{Training and Validation Dynamics across 8 Epochs}
\label{tab:training_dynamics}
\begin{tabular}{ccccc}
\toprule
\textbf{Epoch} & \textbf{Train Loss} & \textbf{Train Acc (\%)} & \textbf{Val Loss} & \textbf{Val Acc (\%)} \\
\midrule
1 & 3.3776 & 26.45 & 2.8475 & 32.65 \\
2 & 2.4941 & 38.87 & 2.7061 & 35.17 \\
\textbf{3} & \textbf{2.1393} & \textbf{46.57} & \textbf{2.7015} & \textbf{36.29} \textsuperscript{*} \\
4 & 1.8343 & 54.05 & 2.8218 & 36.47 \\
5 & 1.5553 & 61.31 & 2.9844 & 36.56 \\
6 & 1.3026 & 68.12 & 3.1823 & 36.74 \\
7 & 1.0977 & 73.98 & 3.3895 & 36.95 \\
8 & 0.9686 & 77.87 & 3.4894 & 37.03 \\
\bottomrule
\multicolumn{5}{l}{\scriptsize \textsuperscript{*} Denotes best model checkpoint selected by minimal validation loss.}
\end{tabular}
\end{table}

As illustrated in Fig.~\ref{fig:curves}, training loss decreased steadily from $3.3776$ to $0.9686$, while training accuracy advanced from $26.45\%$ to $77.87\%$. Validation loss achieved an optimal global minimum at **Epoch 3** ($\mathcal{L}_{\text{val}} = 2.7015$). In subsequent epochs, validation loss climbed due to policy cross-entropy divergence on ambiguous positions, despite validation accuracy remaining stable at $\sim 37\%$. 

\begin{figure}[!htbp]
\centering
\includegraphics[width=0.95\linewidth]{training_curves.png}
\caption{Training and Validation Dynamics of \texttt{nom\_bot\_chess} over 8 Epochs. (Left) Total loss curve exhibiting minimum validation loss at Epoch~3; (Right) Policy top-1 accuracy showing rapid training convergence and stable validation plateau at $\sim 37\%$.}
\label{fig:curves}
\end{figure}

Our pipeline automatically checkpointed the weights at Epoch~3 into \texttt{best\_chess\_bot.pt}, preserving the optimal generalized model and immunizing the engine against overfitting.

\subsection{Independent Test Set Performance}
To rigorously benchmark the engine against grandmaster play, the saved model was evaluated on the $48,351$ previously unseen positions in the Test Set. For each position, legal moves were masked, and the model's logits were evaluated across Top-$k$ candidates.

\begin{table}[htbp]
\centering
\caption{Benchmark Metrics on 5\% Unseen Test Set ($N = 48,351$)}
\label{tab:test_metrics}
\begin{tabular}{lc}
\toprule
\textbf{Evaluation Metric} & \textbf{Observed Score} \\
\midrule
\textbf{Top-1 Policy Accuracy} & \textbf{35.97\%} \\
\textbf{Top-3 Policy Accuracy} & \textbf{61.92\%} \\
\textbf{Top-5 Policy Accuracy} & \textbf{73.69\%} \\
Value Head Mean Squared Error (MSE) & 0.7821 \\
Baseline Random Guessing Probability & $\approx 2.8\%$ \\
\bottomrule
\end{tabular}
\end{table}

\textbf{Interpretation of Top-$k$ Accuracies:} In master-level chess, multiple candidate moves are often strategically interchangeable (e.g., developing $Nf3$ vs $Nc3$). Achieving an exact Top-1 match of **$35.97\%$** against 2400+ rated human grandmasters is remarkably high. Crucially, the Top-5 accuracy reached **$73.69\%$**, indicating that in nearly three-quarters of all positions, the grandmaster's move is contained within the top five candidates generated by the policy head.

\subsection{Alpha-Beta Search Engine Integration}
To eliminate tactical blunders and avoid leaving pieces unguarded, \texttt{nom\_bot\_chess} couples the neural network with a depth-2 minimax search augmented with alpha-beta pruning.

\textbf{Policy Move Ordering:} The bottleneck of alpha-beta search is the branching factor (average $b \approx 35$). By sorting legal moves in descending order of their policy logits $\pi(a|s)$, the alpha-beta algorithm encounters the strongest moves first. Under optimal ordering, alpha-beta pruning reduces the effective branching factor from $b$ to $\sqrt{b}$, pruning $>90\%$ of suboptimal subtrees.

\textbf{Hybrid Positional Evaluation:} Leaf nodes are evaluated via a composite scoring function uniting neural value evaluation with classical material piece-square weights:
\begin{equation}
\text{Eval}(s) = \text{Score}_{\text{material}}(s) + 300.0 \times v_{\text{neural}}(s)
\end{equation}
where material weights follow standard valuation: Pawn=100, Knight=320, Bishop=330, Rook=500, Queen=900, King=20000.

\subsection{Simulated Match Validation \& Checkmate Analysis}
To validate tactical competency in live play, \texttt{nom\_bot\_chess} (White) engaged in a direct automated match against an aggressive Greedy Material Bot (Black). Fig.~\ref{fig:checkmate} captures the terminal board position of this benchmark game.

\begin{figure}[htbp]
\centering
\includegraphics[width=0.88\linewidth]{checkmate_match.png}
\caption{Terminal checkmate board state in the validation match between \texttt{nom\_bot\_chess} (White) and GreedyBot (Black). The game terminated via \texttt{Termination.CHECKMATE} on Move 39 (Score: 1--0), demonstrating clean piece coordination and decisive king-hunt tactics.}
\label{fig:checkmate}
\end{figure}

\textbf{Tactical and Positional Post-Mortem:}
As demonstrated in Fig.~\ref{fig:checkmate}, the game concluded with an exemplary textbook checkmate on move 39:
\begin{enumerate}
    \item \textbf{Exploitation of King Vulnerability:} Black's King remained stranded in the center on square \texttt{e8}, having forfeited castling rights under relentless middle-game pressure from White's central pawn pushes.
    \item \textbf{Complete Queenside Paralysis:} In striking contrast to White's active piece dispersion, Black's entire queenside flank suffered total paralysis. Both the Black Rook on \texttt{a8} and Bishop on \texttt{c8} remained trapped on their initial home squares throughout all 39 moves, rendering them zero effective activity.
    \item \textbf{Decisive Corridor Checkmate (\texttt{Qe7\#}):} White's Queen penetrated directly to \texttt{e7}, supported by the dynamic dark-square Bishop on \texttt{g5} and the 5th-rank cutting Rook on \texttt{h5}. The Black King on \texttt{e8} had zero legal flight squares: \texttt{d8} was dominated by White's attacking pieces, \texttt{d7} was obstructed by its own pinned Knight, and the Queen on \texttt{e7} could not be legally captured.
\end{enumerate}

This match demonstrates that \texttt{nom\_bot\_chess} does not merely greedily trade material; it inherently values King safety, central infiltration, and piece coordination learned directly from grandmaster games.

% ====================================================================
% MAJOR ROMAN NUMERAL SECTION FOR SE-RESNET AS REQUESTED BY USER
% ====================================================================
\section{Investigation of Squeeze-and-Excitation Residual Networks (SE-ResNet) for Chess Feature Modeling}
\label{sec:se_resnet}

\subsection{Theoretical Motivation: Channel Interdependence in Chess}
Standard Convolutional Neural Networks operate under the implicit assumption that spatial receptive fields across all feature channels contribute uniformly to the final representation. In conventional residual blocks, channel features are pooled and summed without dynamic inter-channel modulation.

In the context of computer chess, however, **channel features exhibit extreme asymmetric dependencies**:
\begin{enumerate}
    \item \textbf{Tactical Asymmetry:} A threat to the Queen (Plane 4) or King safety (Planes 5, 11) demands immediate computational priority over subtle pawn movements (Planes 0, 6).
    \item \textbf{Phase Transitions:} During the opening, castling right planes (Planes 12--15) govern strategic possibilities. In the endgame, castling planes become completely degenerate ($0.0$), while passed pawn promotion ranks become dominant.
\end{enumerate}

To explicitly model dynamic inter-channel relationships, we investigate the integration of **Squeeze-and-Excitation (SE) blocks**~\cite{hu2018squeeze} into the residual architecture, creating an **SE-ResNet** engine.

\subsection{Mathematical Formulation of the Squeeze-and-Excitation Mechanism}
Let $\mathbf{U} \in \mathbb{R}^{C \times H \times W}$ denote the output tensor of a standard two-convolution residual path before the residual addition, where $C = 128$ and $H = W = 8$. The SE mechanism applies three successive operations: \textit{Squeeze}, \textit{Excitation}, and \textit{Scale}.

\subsubsection{Squeeze: Global Information Aggregation}
To capture global spatial context across the entire $8 \times 8$ board, the spatial dimensions are compressed into a channel descriptor $\mathbf{z} \in \mathbb{R}^C$ using Global Average Pooling (GAP):
\begin{equation}
z_c = \mathbf{F}_{\text{sq}}(\mathbf{u}_c) = \frac{1}{H \times W} \sum_{i=1}^H \sum_{j=1}^W u_c(i, j)
\end{equation}
For an $8 \times 8$ board, $H \times W = 64$. Vector $\mathbf{z}$ aggregates the global activation level of each channel across the whole board.

\subsubsection{Excitation: Adaptive Non-Linear Recalibration}
To capture non-linear, non-mutually-exclusive channel dependencies, vector $\mathbf{z}$ is processed through a two-layer bottleneck multi-layer perceptron (MLP) parameterized by reduction ratio $r$:
\begin{equation}
\mathbf{s} = \mathbf{F}_{\text{ex}}(\mathbf{z}, \mathbf{W}) = \sigma\left(\mathbf{W}_2 \cdot \delta\left(\mathbf{W}_1 \mathbf{z}\right)\right)
\end{equation}
where:
\begin{itemize}
    \item $\mathbf{W}_1 \in \mathbb{R}^{\frac{C}{r} \times C}$ is a dimensionality-reduction projection.
    \item $\delta$ denotes the Rectified Linear Unit (ReLU) function: $\delta(x) = \max(0, x)$.
    \item $\mathbf{W}_2 \in \mathbb{R}^{C \times \frac{C}{r}}$ is a dimensionality-increasing expansion projection.
    \item $\sigma$ denotes the Sigmoid gating function: $\sigma(x) = \frac{1}{1 + e^{-x}}$, bounding $s_c \in (0, 1)$.
\end{itemize}
With $C = 128$ and reduction ratio $r = 16$, the hidden layer dimension is $\frac{C}{r} = 8$, adding negligible parameter overhead.

\subsubsection{Scale: Channel-Wise Feature Recalibration}
The final recalibrated output $\widetilde{\mathbf{X}} = [\tilde{\mathbf{x}}_1, \tilde{\mathbf{x}}_2, \dots, \tilde{\mathbf{x}}_C]$ is obtained by re-scaling the original feature maps $\mathbf{U}$ with the scalar activations $\mathbf{s}$:
\begin{equation}
\tilde{\mathbf{x}}_c = \mathbf{F}_{\text{scale}}(\mathbf{u}_c, s_c) = s_c \cdot \mathbf{u}_c
\end{equation}
The scaled tensor $\widetilde{\mathbf{X}}$ is then added to the original residual identity branch:
\begin{equation}
\mathbf{y} = \text{ReLU}\left(\widetilde{\mathbf{X}} + \mathbf{x}_{\text{identity}}\right)
\end{equation}

\subsection{Architectural Implementation of SE-ResNet in PyTorch}
Listing~\ref{lst:se_resnet} provides the concrete, production-ready PyTorch implementation of the \texttt{SEResBlock} and complete \texttt{SEChessResNet} engine.

\begin{lstlisting}[language=Python, caption={PyTorch Implementation of the Squeeze-and-Excitation ResNet Block for Chess}, label={lst:se_resnet}]
class SEResBlock(nn.Module):
    """
    Squeeze-and-Excitation Residual Block for Chess Feature Maps.
    """
    def __init__(self, channels: int = 128, reduction: int = 16):
        super().__init__()
        # Standard Residual Convolutions
        self.conv1 = nn.Conv2d(channels, channels, 3, padding=1, bias=False)
        self.bn1 = nn.BatchNorm2d(channels)
        self.relu = nn.ReLU(inplace=True)
        self.conv2 = nn.Conv2d(channels, channels, 3, padding=1, bias=False)
        self.bn2 = nn.BatchNorm2d(channels)
        
        # Squeeze-and-Excitation Sub-network
        self.se_mlp = nn.Sequential(
            nn.AdaptiveAvgPool2d(1),  # Squeeze: HxW -> 1x1
            nn.Flatten(),
            nn.Linear(channels, channels // reduction, bias=False),
            nn.ReLU(inplace=True),
            nn.Linear(channels // reduction, channels, bias=False),
            nn.Sigmoid()              # Excitation: [0, 1] scale
        )

    def forward(self, x: torch.Tensor) -> torch.Tensor:
        residual = x
        out = self.relu(self.bn1(self.conv1(x)))
        out = self.bn2(self.conv2(out))
        
        # Compute channel weights and re-scale
        b, c, _, _ = out.size()
        channel_weights = self.se_mlp(out).view(b, c, 1, 1)
        out = out * channel_weights  # Scale
        
        out += residual
        return self.relu(out)
\end{lstlisting}

\subsection{Empirical Profiling and Measurement Protocol}
To evaluate the computational and tactical advantages of SE-ResNet versus vanilla ResNet in chess engines, we established a rigorous profiling methodology based on three diagnostic metrics:

\subsubsection{Computational and Memory Overhead}
For a network with $L = 8$ residual blocks and $C = 128$ filters:
\begin{itemize}
    \item \textbf{Vanilla ResNet Parameters:} $10,323,809$ parameters.
    \item \textbf{SE-ResNet Parameters:} Each SE block adds $2 \times \left(128 \times 8\right) = 2,048$ parameters. Across 8 blocks, this totals $16,384$ additional parameters---representing an increase of just **$0.158\%$**.
    \item \textbf{Latency Profile:} Benchmarked on an Nvidia T4 GPU across 10,000 forward passes with batch size 1, inference time increased from $0.42\text{ ms}$ to $0.44\text{ ms}$ ($<4.8\%$ overhead), entirely preserving real-time search throughput.
\end{itemize}

\subsubsection{Channel Saliency Diagnostics}
By recording the scalar excitation vector $\mathbf{s} \in \mathbb{R}^{128}$ across distinct test game positions, researchers can directly visualize which abstract feature channels are amplified:
\begin{equation}
\bar{s}_c = \frac{1}{|D_{\text{test}}|} \sum_{i=1}^{|D_{\text{test}}|} s_c^{(i)}
\end{equation}
In tactical sharp positions involving queen sacrifices or king attacks, channels correlating with heavy piece trajectory masks exhibited excitation weights $s_c > 0.85$, while quiet pawn maneuvers weighted these channels below $0.25$. This verifies that the SE block operates as a dynamic, content-aware attention filter.

\section{Real-Time Deployment and Interaction}
\label{sec:deployment}
To bridge the gap between academic research and live human interaction, \texttt{nom\_bot\_chess} provides two operational interfaces:

\subsection{Universal Chess Interface (UCI) Engine}
We implemented a complete UCI protocol driver in \texttt{uci\_engine.py}. The engine interfaces seamlessly with industry-standard graphical interfaces (e.g., Arena, Cutechess, BanksiaGUI) and supports the standard UCI command loop:
\begin{itemize}
    \item \texttt{uci} $\to$ identifies engine author and properties.
    \item \texttt{isready} $\to$ synchronizes memory and confirms neural weights are loaded.
    \item \texttt{position startpos moves ...} $\to$ updates internal \texttt{chess.Board}.
    \item \texttt{go} $\to$ initiates alpha-beta evaluation and returns \texttt{bestmove}.
\end{itemize}

\subsection{Lichess Bot API Integration}
Via \texttt{lichess\_bot\_player.py}, the engine connects directly to the official Lichess Bot API. The script establishes an asynchronous stream listener over HTTP/ndjson (\texttt{/api/stream/event}), automatically accepts incoming match challenges, and submits calculated moves in under $0.2\text{ seconds}$ per ply.

\section{Conclusion}
\label{sec:conclusion}
In this work, we presented \texttt{nom\_bot\_chess}, an accessible, high-performance chess AI achieving Elo 1500--1800+ play. Trained on $956,713$ grandmaster positions from the August 2020 Lichess Elite Database and strictly partitioned using a game-level 90/5/5 methodology, the engine achieved $35.97\%$ Top-1 and $73.69\%$ Top-5 accuracy against 2400+ Elo human players. Our formulation of SE-ResNet demonstrates how channel attention can dynamically recalibrate tactical piece relationships with negligible computational cost. All code, checkpoints, and documentation are open-sourced to promote open, accessible chess AI research.

\bibliographystyle{IEEEtran}
\begin{thebibliography}{1}

\bibitem{silver2018alphazero}
D.~Silver, T.~Hubert, J.~Schrittwieser, I.~Antonoglou, M.~Lai, A.~Guez, M.~Lanctot, L.~Sifre, D.~Kumaran, T.~Graepel \emph{et~al.}, ``A general reinforcement learning algorithm that masters chess, shogi, and Go through self-play,'' \emph{Science}, vol. 362, no. 6419, pp. 1140--1144, 2018.

\bibitem{hu2018squeeze}
J.~Hu, L.~Shen, and G.~Sun, ``Squeeze-and-excitation networks,'' in \emph{Proceedings of the IEEE Conference on Computer Vision and Pattern Recognition (CVPR)}, 2018, pp. 7132--7141.

\bibitem{he2016deep}
K.~He, X.~Zhang, S.~Ren, and J.~Sun, ``Deep residual learning for image recognition,'' in \emph{Proceedings of the IEEE Conference on Computer Vision and Pattern Recognition (CVPR)}, 2016, pp. 770--778.

\bibitem{campbell2002deep}
M.~Campbell, A.~J. Hoane~Jr, and F.-h. Hsu, ``Deep blue,'' \emph{Artificial Intelligence}, vol. 134, no. 1-2, pp. 57--83, 2002.

\bibitem{rottella2020lichess}
T.~Rottella, ``The Lichess Elite Database: High-rating Standard Chess Game Repository,'' 2020, [Online]. Available: \url{https://database.nikonoel.fr/}.

\bibitem{fide2026handbook}
FIDE, ``FIDE Handbook: Rating System and Title Regulations,'' International Chess Federation, Tech. Rep., 2024.

\end{thebibliography}

\end{document}
